\subsection{Die Rolle der Industrie im Faschismus}

Die historische Erfahrung Italiens und Deutschlands zeigt: Der Faschismus ist nicht aus dem Nichts entstanden – er ist das Ergebnis einer bewussten, strategischen Einmischung der Industrie in die Politik.

\subsubsection{Italien: Die Industriellen finanzieren Mussolini}

Nach dem Ersten Weltkrieg war Italien in einer tiefen Krise. Die Arbeiterbewegung war stark – 1919 wurde die Sozialistische Partei mit 32 Prozent stärkste Kraft. Es kam zu Fabrikbesetzungen, Landbesetzungen, Streiks. Die Industriellen und Großgrundbesitzer fürchteten eine sozialistische Revolution.

\begin{itemize}
    \item Die \textbf{Confindustria} (der italienische Industriellenverband) finanzierte Mussolinis \emph{Fasci di Combattimento} ab 1920 massiv.
    \item Der \textbf{March auf Rom} 1922 wurde von den Eliten nicht verhindert, sondern ermöglicht – König Vittorio Emanuele III. weigerte sich, den Ausnahmezustand auszurufen.
    \item Die \textbf{Zerschlagung der Gewerkschaften} 1925/26 war die erste Amtshandlung der faschistischen Regierung – und sie wurde von der Industrie begeistert aufgenommen.
\end{itemize}

Die Industriellen glaubten, Mussolini kontrollieren zu können. Sie nutzten ihn als Werkzeug gegen die Arbeiterbewegung – und wurden später selbst von ihm gefressen. Aber der Mechanismus war klar: \emph{Die Industrie rief den Faschismus, um die Demokratie zu zerstören – weil die Demokratie den Arbeitern eine Mehrheit hätte geben können.}

\subsubsection{Deutschland: Die Industriellen finanzieren Hitler}

Das gleiche Muster wiederholte sich in Deutschland:

\begin{itemize}
    \item Die \textbf{Industriellen} (Thyssen, Krupp, IG Farben, Deutsche Bank, Bosch) spendeten ab 1930 massiv an die NSDAP.
    \item Das \textbf{Harzburger Front}-Treffen 1931 vereinte NSDAP, DNVP und Stahlhelm – finanziert von der Industrie.
    \item Die \textbf{Ernennung Hitlers zum Reichskanzler} am 30. Januar 1933 geschah nicht, weil Hitler die Mehrheit hatte (er hatte 33,1 Prozent), sondern weil die konservativen Eliten (Papen, Hugenberg, Hindenburg) ihn als Werkzeug gegen die Arbeiterbewegung einsetzen wollten.
    \item Die \textbf{Reichstagsbrandverordnung} und das \textbf{Ermächtigungsgesetz} wurden von den bürgerlichen Parteien mitgetragen – nicht weil sie Hitler liebten, sondern weil sie die KPD und SPD zerschlagen wollten.
\end{itemize}

Die Industriellen glaubten, Hitler kontrollieren zu können. Sie nutzten ihn als Werkzeug gegen die Arbeiterbewegung – und wurden später selbst von ihm gefressen. Aber der Mechanismus war derselbe: \emph{Die Industrie rief den Faschismus, um die Demokratie zu zerstören – weil die Demokratie den Arbeitern eine Mehrheit hätte geben können.}

\subsubsection{Der Holocaust als industrielles Projekt}

Der Holocaust war nicht nur ein Verbrechen aus Judenhass – er war ein \emph{industrielles Projekt}. Die Vernichtung der Juden wurde möglich durch:

\begin{itemize}
    \item \textbf{Technologie:} Die Eisenbahnen der Deutschen Reichsbahn, die Gaskammern, die Krematorien – alles Industrieprodukte, die ohne die Zusammenarbeit mit der Industrie nicht möglich gewesen wären.
    \item \textbf{Organisation:} Die Bürokratie des Holocaust – Deportationslisten, Transportpläne, Todesstatistiken – war eine logistische Meisterleistung der industriellen Verwaltung.
    \item \textbf{Kapital:} Die Vernichtung war teuer – sie wurde von der Industrie mitfinanziert, die im Gegenzug die Zwangsarbeiter der Konzentrationslager erhielt.
    \item \textbf{Beteiligung der Industrie:} IG Farben, Siemens, BMW, Krupp, Deutsche Bank – alle profitierten von der Zwangsarbeit und der Vernichtung der Juden.
\end{itemize}

Der Holocaust war kein Betriebsunfall der Geschichte. Er war das \emph{logische Ergebnis} einer Verbindung von Judenhass und industrieller Kapitalmacht. Ohne die Industrie – ohne ihre Technologie, ihre Organisation, ihr Kapital – wäre der Holocaust nicht möglich gewesen. Die Verbindung von Judenhass und industrieller Kapitalmacht ist das, was den Holocaust ermöglichte.

\subsubsection{Die strukturelle Kontinuität}

Die Parallele zwischen Italien und Deutschland ist eindeutig:

\begin{table}[h!]
\centering
\begin{tabular}{|l|l|p{5.5cm}|}
\hline
\textbf{Element} & \textbf{Italien} & \textbf{Deutschland} \\
\hline
Bedrohungswahrnehmung & Biennio Rosso (1919–1920) & Novemberrevolution (1918–1919) \\
\hline
Finanzierung & Confindustria, Großgrundbesitzer & Thyssen, Krupp, IG Farben, Deutsche Bank \\
\hline
Politisches Werkzeug & Mussolini & Hitler \\
\hline
Ziel & Zerschlagung der Arbeiterbewegung & Zerschlagung der Arbeiterbewegung \\
\hline
Ergebnis & Faschistische Diktatur & Nationalsozialistische Diktatur \\
\hline
\end{tabular}
\caption{Die strukturelle Parallele zwischen Italien und Deutschland}
\end{table}

Die entscheidende Einsicht ist: \emph{Der Faschismus ist die politische Form, die das Kapital annimmt, wenn die Demokratie ihm gefährlich wird.} Er ist das Ergebnis der industriellen Einmischung in die Politik – und er ist immer antidemokratisch.

Diese Erkenntnis ist nicht nur historisch. Sie ist der Schlüssel zur Gegenwart.